\chapter*{Előszó}
\addcontentsline{toc}{chapter}{Előszó}

Amióta az eszemet tudom, a reáltudományok és a művészetek -- különösen a zene -- világa meghatározó szerepet játszik az életemben. 
Időről időre az egyik kerül előtérbe, majd a másik veszi át a szerepet, de sosem maradt ki teljesen egyik sem az érdeklődésemből.
Ez az oka, hogy sosem tudtam csak egy oldal mellett elköteleződni véglegesen és minden fókuszomat kizárólag ráhelyezni, hogy minél jobb legyek benne. 
Talán nincs is rá szükség. A reáltudományok segítenek leírni a körülöttem lévő fizikai valóságot és lépést tarthatok vele a technológia exponenciális előretörésével. 
Ennek ellenére mégis hiányérzetem támad, ha a saját belső világom szeretném megfigyelni. A művészetek körülírják és teljesebbé teszik az ember lelkivilágát, 
így teremtve meg számomra az egyensúlyt.

Ez az egyensúlykeresés inspirálta jelen kutatásomat és diplomadolgozatomat is, amely a zene és a számítástechnika határterületén valósul meg.

\medskip

\begin{flushright}
  \textit{Nagy Csaba}
\end{flushright}
